% ! TEX root = ../ag-20-21.tex

\chapter{Recapitulare: Metoda lui Gauss}

\section{Eliminare gaussiană}
\index{eliminare gaussiană}

Metoda eliminării gaussiene (numită și metoda Gauss(-Jordan)) este folosită
pentru a aduce matrice la o formă mai simplă, anume triunghiulară sau chiar
forma matricei unitate.

Principalele aplicații ale eliminării gaussiene sînt în rezolvarea sistemelor
de ecuații liniare și în inversarea matricelor.

În ambele situații, se trece de la o stare la alta făcînd \emph{transformări %
  elementare}, adică acele operații permise în determinanți, care nu schimbă
valoarea determinantului. În general, operațiile pe care le putem face intră
sub denumirea de \emph{combinații liniare cu linii sau coloane}.

Un exemplu este următorul. Fie sistemul de ecuații:
\[
  \begin{cases}
    mx_1 + x_2 + x_3 &= 1 \\
    x_1 + mx_2 + x_3 &= m \\
    x_1 + x_2 + mx_3 &= m^2,
  \end{cases}
\]
unde $ m \in \RR $. Evident, soluția sistemului va implica o discuție după $ m $.

Fie $ A $ matricea sistemului, iar $ B $ matricea-coloană a termenilor liberi.
Prelucrăm matricea extinsă $ (A \mid B) $ pînă cînd aducem pe $ A $ în formă
(superior) triunghiulară.

Considerăm matricea $M = (A \mid B) $ de mai sus:
\begin{align*}
  (A | B) &= \begin{pmatrix}
    m & 1 & 1 & \mid & 1 \\
    1 & m & 1 & \mid & m \\
    1 & 1 & m & \mid & m^2
  \end{pmatrix} \\
          &\xrightarrow{L_1 \leftrightarrow L_2}
            \begin{pmatrix}
              1 & m & 1 & \mid & m \\
              m & 1 & 1 & \mid & 1 \\
              1 & 1 & m & \mid & m^2
            \end{pmatrix} \\
          &\stackrel{L_2 \to L_2 - mL_1}{\xrightarrow{L_3 \to L_3 - L_1}}
            \begin{pmatrix}
              1 & m & 1 & \mid & m \\
              0 & 1 - m^2 & 1 - m & \mid & 1 - m^2 \\
              0 & 1 - m & m - 1 & \mid & m^2 - m.
            \end{pmatrix}
\end{align*}

În acest punct, avem o discuție:

(a) Dacă $m = 1$, atunci sistemul se reduce la prima ecuație, deci este compatibil
dublu nedeterminat. Rezultă soluția
\[
  \{ (1 - \alpha - \beta, \alpha, \beta) \mid \alpha, \beta \in \RR \}.
\]

(b) Dacă $m \neq 1$, atunci putem continua transformările și ajungem, în fine, la:
\[
  M = \begin{pmatrix}
    1 & 0 & 1 + m & \mid & m + m^2 \\
    0 & 1 & -1 & \mid & -m \\
    0 & 0 & 2 & \mid & (m + 1)^2
  \end{pmatrix}
\]

Aici discutăm din nou:

(b1) Dacă $m = -2$, sistemul este incompatibil, deoarece ultima ecuație devine $0 = 1$.

(b2) Dacă $m \neq -2$, putem continua transformările și ajungem, în cele din urmă, la:
\[
  M = %
  \begin{pmatrix}
    1 & 0 & 0 & \mid & -(m+1)/2 \\
    0 & 1 & 0 & \mid & 1/(m+2) \\
    0 & 0 & 1 & \mid & (m+1)^2/(m+2)
  \end{pmatrix}
\]


În acest ultim caz, sistemul este compatibil determinat, soluția fiind dată de ultima coloană:
\[
  \begin{cases}
    x_1  &= -\dfrac{m + 1}{2} \\
    & \\
    x_2  &= \dfrac{1}{m + 2} \\
    & \\
    x_3  &= \dfrac{(m + 1)^2}{m + 2}.
  \end{cases}
\]

Concluzia generală este:
\begin{enumerate}[(a)]
\item Dacă $m \in \RR - \{-2, 1\}$, sistemul are soluție unică;
\item Dacă $m = -2$, sistemul este incompatibil;
\item Dacă $m = -1$, sistemul este compatibil nedeterminat.
\end{enumerate}

\vspace{1cm}

Similar, pentru a calcula inversa unei matrice, bordăm matricea dată cu matricea
unitate și efectuăm transformări elementare pînă ce matricea inițială devine
matricea unitate.

Iată un exemplu, notînd și transformările:
\begin{align*}
  \begin{pmatrix}
    2 & 1 & 3 & \mid & 1 & 0 & 0 \\
    -2 & 3 & 4 & \mid & 0 & 1 & 0 \\
    5 & 1 & 1 & \mid & 0 & 0 & 1
  \end{pmatrix} &\xrightarrow{L_1 \to (1/2)L_1}
                  \begin{pmatrix}
                    1 & 1/2 & 3/2 & \mid & 1/2 & 0 & 0 \\
                    -2 & 3 & 4 & \mid & 0 & 1 & 0 \\
                    5 & 1 & 1 & \mid & 0 & 0 & 1 
                  \end{pmatrix} \\
      &\stackrel{L_2 \to 2L_1 + L_2}{\xrightarrow{L_3 \to -5L_1 + L_3}}
        \begin{pmatrix}
          1 & 1/2 & 3/2 & \mid & 1/2 & 0 & 0 \\
          0 & 4 & 7 & \mid & 1 & 1 & 0 \\
          0 & -3/2 & -13/2 & \mid & -5/2 & 0 & 1
        \end{pmatrix} \\
      &\xrightarrow{L_2 \to (1/4)L_2}
        \begin{pmatrix}
          1 & 1/2 & 3/2 & \mid & 1/2 & 0 & 0 \\
          0 & 1 & 7/4 & \mid & 1/4 & 1/4 & 0 \\
          0 & -3/2 & -13/2 & \mid & -5/2 & 0 & 1 
        \end{pmatrix} \\
      &\stackrel{L_1 \to (-1/2)L_2 + L_1}{\xrightarrow{L_3 \to (3/2)L_2 + L_3}}
        \begin{pmatrix}
          1 & 0 & 5/8 & \mid & 3/8 & -1/8 & 0 \\
          0 & 1 & 7/4 & \mid & 1/4 & 1/4 & 0 \\
          0 & 0 & -31/8 & \mid & -17/8 & 3/8 & 1
        \end{pmatrix} \\
      &\xrightarrow{L_3 \to (-8/31)L_3}
        \begin{pmatrix}
          1 & 0 & 5/8 & \mid & 3/8 & -1/8 & 0 \\
          0 & 1 & 7/4 & \mid & 1/4 & 1/4 & 0 \\
          0 & 0 & 1 & \mid & 17/31 & -3/31 & -8/31
        \end{pmatrix} \\
      &\stackrel{L_1 \to (-5/8)L_3 + L_1}{\xrightarrow{L_2 \to (-7/4)L_3 + L_2}}
        \begin{pmatrix}
          1 & 0 & 0 & \mid & 1/31 & -2/31 & 5/31 \\
          0 & 1 & 0 & \mid & -22/31 & 13/31 & 14/31 \\
          0 & 0 & 1 & \mid & 17/31 & -3/31 & -8/31
        \end{pmatrix}
\end{align*}

Rezultă că 
\[
  A^{-1} = %
  \begin{pmatrix} 
    1/31 & -2/31 & 5/31 \\
    -22/31 & 13/31 & 14/31 \\
    17/31 & -3/31 & -8/31
  \end{pmatrix}.
\]

%%%%%%%%%%%%%%%%%%%%%%%%%%%%%%%%%%%%%%%%%%%%%%%%%%%%%%%%%%%%%%%%%%%%%% 

\subsection{Exerciții}

1. Rezolvați următoarele sisteme, atît cu metoda matriceală clasică, cît și cu
metoda lui Gauss:
\begin{enumerate}[(a)]
\item $ \begin{cases}
    x + 3y - z &= 1 \\
    3x + y + 2z &= 4 \\
    5x - 2y + z &= 2
  \end{cases} $
\item $ \begin{cases}
    3x - y + 2z + 3t &= 4 \\
    x + 3y - z + t &= -1
  \end{cases} $
\item $ \begin{cases}
    2x - y + 3z &= 1 \\
    3x + 2y - z &= 4 \\
    x + 3y - 4z &= 3
  \end{cases} $.
\end{enumerate}

\vspace{1cm}

2. Calculați inversele matricelor, atît cu metoda folosind matricea adjunctă,
cît și folosind metoda lui Gauss:
\begin{enumerate}[(a)]
\item $ A =
  \begin{pmatrix}
    1 & 0 & -1 \\
    3 & 2 & 1 \\
    0 & 1 & 1
  \end{pmatrix} $;
\item $ B =
  \begin{pmatrix}
    -1 & 2 & 4 & 1 \\
    3 & 0 & 1 & 4 \\
    0 & 1 & 1 & 0 \\
    2 & 1 & 0 & 1
  \end{pmatrix} $;
\item $ C =
  \begin{pmatrix}
    1 & 1 & 2 \\
    0 & 1 & 0 \\
    0 & 0 & -1
  \end{pmatrix} $.
\end{enumerate}


%%% Local Variables:
%%% mode: latex
%%% TeX-master: "../ag-20-21"
%%% End:
